\documentclass{ieeeaccess}
\usepackage{cite}
\usepackage{amsmath,amssymb,amsfonts}
\usepackage{algorithmic}
\usepackage{graphicx}
\usepackage{textcomp}
\usepackage{booktabs}
\usepackage{multirow}
\usepackage{url}
\usepackage{algorithm}

\def\BibTeX{{\rm B\kern-.05em{\sc i\kern-.025em b}\kern-.08em
    T\kern-.1667em\lower.7ex\hbox{E}\kern-.125emX}}

\begin{document}
\history{Date of publication October 00, 2026, date of current version October 00, 2026.}
\doi{10.1109/ACCESS.2026.DOI}

\title{NeuroClean: High-Fidelity Smartphone Image Denoising via Hybrid U-Net--DnCNN Dual-Stream Deep Learning and Overlapping Hann-Window Tiling}

\author{\uppercase{Majji Naveen}\authorrefmark{1},
\uppercase{Parri Hari Chandan}\authorrefmark{1},
\uppercase{M. Siva Charan Reddy}\authorrefmark{1},
\uppercase{and Jayalakshmi M}\authorrefmark{1*}, \IEEEmembership{Senior Member, IEEE}}

\address[1]{School of Computer Science and Engineering (SCOPE), Vellore Institute of Technology (VIT), Vellore, Tamil Nadu 632014, India}
\tfootnote{This work was supervised by Dr. Jayalakshmi M at the School of Computer Science and Engineering (SCOPE), Vellore Institute of Technology, Vellore.}

\markboth
{Naveen \headeretal: NeuroClean: High-Fidelity Smartphone Image Denoising via Hybrid U-Net--DnCNN}
{Naveen \headeretal: NeuroClean: High-Fidelity Smartphone Image Denoising via Hybrid U-Net--DnCNN}

\corresp{Corresponding author: Jayalakshmi M (e-mail: m.jayalakshmi@vit.ac.in).}

\begin{abstract}
Smartphone cameras operate under acute physical form-factor limitations: micro-apertures and sub-micron to micron pixel pitches ($1.0\,\mu\text{m} - 1.4\,\mu\text{m}$) restrict photon capture, yielding severe Poisson photon shot noise, amplifier read noise, and chromatic grain under low-ambient illumination. Existing neural restoration approaches either suffer from oversmoothing of high-frequency photographic textures (in deep multi-scale autoencoders like U-Net) or fail to preserve low-frequency luminance continuity across large receptive fields (in feed-forward residual networks like DnCNN). Furthermore, deploying deep models on arbitrary, high-resolution smartphone captures ($12+$ to $100+$ megapixels) induces fatal memory exhaustion or catastrophic boundary grid artifacts under naive spatial tiling. In this paper, we present NeuroClean, a unified deep learning and production deployment framework for smartphone image restoration. NeuroClean introduces a dual-stream Hybrid U-Net--DnCNN neural architecture that synergizes: (i) a multi-scale contracting-expansive U-Net branch with long skip connections to capture scene context and uniform illumination gradients, and (ii) a non-pooling 10-layer residual DnCNN branch operating at full spatial resolution to isolate high-frequency sensor noise without spatial decimation. The feature representations from both streams are fused through a $1 \times 1$ cross-channel projection layer to estimate the underlying residual noise field $\mathcal{R}(y)$, yielding the restored image via direct residual subtraction $\hat{x} = \text{clip}(y - \mathcal{R}_{\text{hybrid}}(y), 0, 1)$. To process arbitrary image dimensions without memory bloat or edge stitching seams, we engineer a 2D Hann-window overlapping cosine tiling engine with reflect padding and weighted overlap-add reconstruction. Evaluated on paired real-world smartphone photographs ($1404 \times 2040$ native resolution), our hybrid model achieves a mean Peak Signal-to-Noise Ratio (PSNR) of $28.60+\,\text{dB}$ and Structural Similarity Index (SSIM) of $0.8920+$, representing a $+4.56\,\text{dB}$ and $+0.1915$ improvement over the raw degraded baseline ($24.04\,\text{dB}$, $0.7005$ SSIM), outperforming standard DnCNN ($27.42\,\text{dB}$, $0.8650$) and 3-level U-Net ($27.95\,\text{dB}$, $0.8710$) with statistically significant gains under Holm-corrected Wilcoxon signed-rank tests ($p < 0.001$). The complete system is deployed as an open-source, production-grade cloud application featuring an asynchronous FastAPI inference engine, a responsive glassmorphic frontend with interactive comparison sliders, real-time blind noise estimation ($\hat{\sigma}$), and an interactive 3D WebGL photodiode sensor simulation.
\end{abstract}

\begin{keywords}
Smartphone image denoising, computational photography, deep residual learning, U-Net, DnCNN, Hann-window tiling, photon shot noise, image sensor degradation, deep learning deployment.
\end{keywords}

\titlepgskip<0.1in>

\maketitle

\section{Introduction}
\label{sec:introduction}
\PARstart{C}{omputational} photography has emerged as the definitive driver of image quality improvements in modern mobile consumer electronics. Over the past decade, smartphones have superseded dedicated digital single-lens reflex (DSLR) and mirrorless cameras for everyday photographic capture \cite{delbracio2021mobile}. However, smartphone camera hardware remains constrained by severe physical limitations. While full-frame and APS-C camera sensors possess pixel pitches between $3.7\,\mu\text{m}$ and $8.4\,\mu\text{m}$, mobile camera sensors must compress tens of millions of photosensitive diodes onto sensor dies smaller than $1/1.3$ inches, yielding pixel pitches typically spanning $0.8\,\mu\text{m}$ to $1.4\,\mu\text{m}$ \cite{barron2022mip}.

Under low-light ambient conditions, high shutter speeds, or high ISO settings, the reduced physical collection area of each photodiode severely restricts the number of incident photons \cite{hasinoff2016burst}. The resulting captures are degraded by composite stochastic noise processes:
\begin{enumerate}
    \item \textbf{Photon shot noise}, arising from the Poisson distribution of photon arrival times \cite{milanfar2013tour};
    \item \textbf{Thermal and analog read noise}, generated by charge-to-voltage conversion amplifiers during analog-to-digital conversion (ADC) \cite{healey1994radiometric};
    \item \textbf{Chromatic noise}, resulting from spatial interpolation during Bayer color filter array (CFA) demosaicing \cite{ramanath2002demosaicking}.
\end{enumerate}

Historically, image denoising relied upon handcrafted spatial and transform-domain filters, such as Bilateral Filtering \cite{tomasi1998bilateral}, Non-Local Means (NLM) \cite{buades2005non}, and Block-Matching and 3D Filtering (BM3D) \cite{dabov2007image}. While mathematically elegant, these techniques assume additive white Gaussian noise (AWGN) with stationary variance across the entire image plane. In real smartphone image signal processors (ISPs), however, noise is signal-dependent, heteroscedastic, and spatially correlated across RGB color channels \cite{foi2008practical}. Consequently, classical filters frequently blur fine organic textures (e.g., foliage, facial hair, fabric weaves) while leaving low-frequency chroma splotches intact.

The advent of deep convolutional neural networks (CNNs) revolutionized image restoration. Zhang \textit{et al.} introduced DnCNN (Denoising Convolutional Neural Network) \cite{zhang2017beyond}, demonstrating that predicting the residual noise field rather than the clean image directly accelerates training convergence and boosts restoration fidelity. Concurrently, Ronneberger \textit{et al.}'s U-Net \cite{ronneberger2015u} architecture---originally developed for biomedical semantic segmentation---gained widespread adoption in image-to-image translation tasks due to its multi-scale contracting-expansive paths and long skip connections that preserve macro-structural context \cite{isola2017image}.

Despite these advances, three critical bottlenecks persist in contemporary smartphone denoising research:
\begin{enumerate}
    \item \textbf{The Context vs. Grain Dilemma:} Feed-forward residual chains (such as standard DnCNN) operate at native image resolution and excel at isolating high-frequency pixel grain, but their effective receptive field is mathematically bounded by network depth ($21 \times 21$ pixels for 10 layers), making them blind to large-scale luminance gradients and uniform background splotches. Conversely, multi-scale encoder-decoder networks (such as U-Net) capture expansive contextual semantics via successive pooling and downsampling, but downsampling decodes fine photographic edges and introduces spatial phase distortions upon upsampling.
    \item \textbf{The High-Resolution Boundary Artifact Barrier:} Smartphone photographs regularly exceed $12$ to $48$ megapixels ($3000 \times 4000$ to $6000 \times 8000$ pixels). Feeding such high-resolution tensors directly into deep neural networks causes catastrophic out-of-memory (OOM) GPU allocation faults. Conversely, dividing the image into naive non-overlapping rectangular blocks yields severe grid-like boundary stitching discontinuities, as receptive field edges suffer from boundary-truncation errors \cite{huang2021evaluation}.
    \item \textbf{The Prototype-to-Deployment Disconnect:} Academic denoising research frequently terminates with isolated MATLAB or Jupyter notebook scripts evaluated solely on downsampled synthetic patches. Real-world end users require an accessible, zero-friction, cross-platform web environment capable of processing user-uploaded images of arbitrary aspect ratios, providing interactive side-by-side split comparisons, calculating live latency and noise variance metrics, and enforcing strict data privacy without permanent server-side retention.
\end{enumerate}

\subsection{Contributions}
To overcome these challenges, this paper presents \textbf{NeuroClean}, an end-to-end deep learning framework and production deployment system for smartphone photograph denoising. Specifically, our contributions are as follows:
\begin{enumerate}
    \item \textbf{Dual-Stream Hybrid Architecture:} We design and implement a novel \textbf{Hybrid U-Net--DnCNN} deep neural network that combines multi-scale contextual feature extraction (Stream 1) with non-pooling high-frequency residual grain modeling (Stream 2) joined via a $1 \times 1$ cross-channel fusion projection layer.
    \item \textbf{Deep Residual Formulation:} We formulate the restoration objective as deep residual subtraction $\hat{x} = \text{clip}(y - \mathcal{R}_{\text{hybrid}}(y), 0, 1)$, proving empirically that learning the noise manifold $\mathcal{R}(y)$ provides superior gradient flow and faster empirical convergence compared to direct clean-image reconstruction.
    \item \textbf{2D Hann-Window Overlapping Cosine Tiling Engine:} We engineer a seamless patch-processing pipeline utilizing 2D separable Hann cosine windows with $50\%$ overlap and reflect-padding. This guarantees zero boundary seam discontinuities and mathematically uniform reconstructed luminance across arbitrary input dimensions without GPU memory overflow.
    \item \textbf{Comprehensive Empirical Benchmarking:} We conduct rigorous quantitative evaluations across real-world smartphone photograph test sets ($1404 \times 2040$ native dimensions), demonstrating that our hybrid model achieves \textbf{$28.60+\,\text{dB}$ PSNR} and \textbf{$0.8920+\,\text{SSIM}$}, outperforming 3-layer baseline CNNs ($26.87\,\text{dB}$), standard DnCNN ($27.42\,\text{dB}$), and standard U-Net ($27.95\,\text{dB}$) with statistically validated significance under Holm-corrected Wilcoxon signed-rank tests ($p < 0.001$).
    \item \textbf{Full-Stack Production Deployment \& 3D WebGL Visualization:} We build and deploy a complete, production-grade cloud web application comprising a high-throughput asynchronous FastAPI backend, a responsive React/TypeScript frontend with glassmorphism UI/UX, Dark/Light theme switching, interactive comparison sliders, and an interactive 3D Three.js WebGL simulation illustrating smartphone photodiode noise mechanics and neural manifold reconstruction.
\end{enumerate}

\subsection{Research Questions}
This paper investigates five specific research questions:
\begin{itemize}
    \item \textbf{RQ1:} How does the proposed Hybrid U-Net--DnCNN dual-stream architecture compare quantitatively and qualitatively against classical filters, baseline CNNs, standard DnCNN, and standard U-Net on real smartphone noise?
    \item \textbf{RQ2:} Does learning the residual noise field $\mathcal{R}(y)$ provide superior restoration fidelity and training stability compared to direct clean image synthesis $\mathcal{G}(y) \approx x$?
    \item \textbf{RQ3:} Does the 2D Hann-window overlapping cosine tiling engine eliminate boundary stitching artifacts without degrading global reconstruction metrics across arbitrary dimensions?
    \item \textbf{RQ4:} What is the empirical latency, memory footprint, and throughput of the hybrid model across heterogeneous hardware environments (CPU vs. GPU)?
    \item \textbf{RQ5:} How effectively does the end-to-end full-stack web application deliver reliable, privacy-preserving, and intuitive restoration for arbitrary real-world smartphone photographs?
\end{itemize}

\section{Related Work}
\label{sec:related_work}

\subsection{Smartphone Image Noise and Sensor Physics}
Unlike consumer DSLR cameras equipped with large physical sensors, smartphone image sensors have photodiodes whose active area typically measures $1.0\,\mu\text{m} \times 1.0\,\mu\text{m}$ to $1.4\,\mu\text{m} \times 1.4\,\mu\text{m}$ \cite{barron2022mip, abdelhamed2018high}. Under photon-starved captures, quantum fluctuations in photon arrival rate dictate that the measured signal follows a Poisson distribution:
\begin{equation}
P(k; \lambda) = \frac{\lambda^k e^{-\lambda}}{k!}
\end{equation}
where $\lambda$ represents the true irradiance \cite{milanfar2013tour}. When combined with thermal agitation (Johnson-Nyquist noise) and amplifier readout fluctuations modeled as zero-mean Gaussian noise $\mathcal{N}(0, \sigma_r^2)$, the composite noise field $v$ is heteroscedastic:
\begin{equation}
\text{Var}(v \mid x) = \alpha x + \sigma_r^2
\end{equation}
where $\alpha$ scales with sensor analog gain (ISO setting) \cite{foi2008practical}. Subsequent ISP stages---including dark-current subtraction, demosaicing, white balancing, and tone-curve mapping---spatially correlate this noise across neighboring pixels, rendering simplistic AWGN assumptions ineffective \cite{plotz2017benchmarking}.

\subsection{Classical Spatial and Transform Filters}
Early computational image restoration relied on local and non-local spatial averaging. Tomasi and Manduchi developed the Bilateral Filter \cite{tomasi1998bilateral}, which weights neighboring pixels according to spatial proximity and radiometric photometric distance. Buades \textit{et al.} introduced Non-Local Means (NLM) \cite{buades2005non}, exploiting self-similarity across non-local patches. Dabov \textit{et al.} established the benchmark standard with BM3D \cite{dabov2007image}, grouping similar 2D image fragments into 3D arrays and performing collaborative transform-domain hard thresholding and Wiener filtering. While effective for moderate synthetic Gaussian noise, BM3D is computationally intensive ($\mathcal{O}(N^2)$ patch search complexity) and introduces visible high-frequency ringing and watercolor-like smearing on real smartphone noise patterns \cite{lebrun2012secrets}.

\subsection{Deep Residual Learning and DnCNN}
Zhang \textit{et al.} pioneered the modern deep learning paradigm for image restoration with DnCNN \cite{zhang2017beyond}. DnCNN incorporates 17 to 20 convolutional layers with Batch Normalization (BN) and Rectified Linear Units (ReLU), explicitly adopting a residual learning formulation:
\begin{equation}
\mathcal{L}(\Theta) = \frac{1}{2N} \sum_{i=1}^N \| \mathcal{R}(y_i; \Theta) - (y_i - x_i) \|_F^2
\end{equation}
DnCNN demonstrated that predicting the residual noise field rather than the clean image directly accelerates training convergence and acts as an implicit regularizer. Subsequent architectures like FFDNet \cite{zhang2018ffdnet} and CBDNet \cite{guo2019toward} incorporated noise estimation sub-networks and asymmetric loss functions. However, because DnCNN maintains identical spatial resolution across all layers without downsampling, its receptive field is strictly limited to $(2D + 1) \times (2D + 1)$ where $D$ is network depth. Consequently, standard DnCNN cannot resolve low-frequency chroma mottling spanning tens or hundreds of pixels.

\subsection{Multi-Scale Contracting-Expansive Architectures and U-Net}
To capture expansive receptive fields without parameter explosion, Ronneberger \textit{et al.} introduced the U-Net architecture \cite{ronneberger2015u}. U-Net uses successive convolutional and max-pooling layers in an encoder branch to downsample spatial dimensions, capturing high-level contextual semantics, followed by a symmetric decoder branch with upsampling convolutions to restore native spatial resolution. High-resolution spatial features from the contracting path are transferred directly to the expansive path via concatenated skip connections. U-Net and its derivatives (e.g., Attention U-Net \cite{oktay2018attention}, Residual U-Net \cite{zhang2018road}) have achieved remarkable success in computational imaging. Nevertheless, the spatial downsampling inherent to U-Net's pooling layers inevitably destroys ultra-fine photographic micro-textures, resulting in noticeable edge softening when restoring complex natural textures.

\subsection{Full-Resolution Image Tiling and Boundary Discontinuity}
Deploying deep convolutional networks to full-resolution photographs requires dividing massive tensors into computationally tractable sub-patches. Naive tiling extracts non-overlapping square blocks, processes each block independently, and stitches them together \cite{huang2021evaluation}. However, because convolution operations at tile boundaries rely on zero-padding or edge replication, receptive fields near the boundary lack exterior context, creating severe orthogonal boundary grid lines across the restored image \cite{szeliski2022computer}. Overlap-add (OLA) methods with linear feathering reduce boundary seams but suffer from non-uniform energy summation. In this work, we demonstrate that a 2D separable Hann cosine window guarantees constant partition of unity under $50\%$ overlap, completely eliminating boundary discontinuities.

\section{Noise Formulation and Mathematical Foundations}
\label{sec:formulation}

\subsection{Observation Model}
Let $x \in \mathbb{R}^{H \times W \times 3}$ denote the latent, noise-free RGB photograph defined over the normalized intensity range $[0, 1]$. Let $y \in \mathbb{R}^{H \times W \times 3}$ represent the observed, noise-corrupted photograph acquired by the smartphone camera sensor:
\begin{equation}
\mathbf{y} = \mathbf{x} + \mathbf{v}
\end{equation}
where $\mathbf{v} \in \mathbb{R}^{H \times W \times 3}$ denotes the composite noise field comprising heteroscedastic signal-dependent Poisson shot noise, zero-mean Gaussian read noise, and demosaicing chromatic perturbations.

In direct image reconstruction, the network is trained to synthesize the latent image:
\begin{equation}
\mathcal{F}_{\text{direct}}(y; \Theta) \approx x
\end{equation}
In deep residual learning, the network is trained to predict exclusively the residual noise field:
\begin{equation}
\mathcal{R}(y; \Theta) \approx v = y - x
\end{equation}
The reconstructed clean estimate $\hat{x}$ is subsequently recovered via residual subtraction and hard range clipping:
\begin{equation}
\hat{x} = \text{clip}\left( y - \mathcal{R}(y; \Theta), \, 0.0, \, 1.0 \right)
\end{equation}

\subsection{Image Quality Metrics}
Given a ground-truth reference image $x$ and a restored estimate $\hat{x}$ of size $H \times W$ with $C=3$ color channels, the Mean Squared Error (MSE) is:
\begin{equation}
\text{MSE}(x, \hat{x}) = \frac{1}{H W C} \sum_{i=1}^H \sum_{j=1}^W \sum_{c=1}^C \left( x(i, j, c) - \hat{x}(i, j, c) \right)^2
\end{equation}
The Peak Signal-to-Noise Ratio (PSNR) expressed in decibels (dB) is:
\begin{equation}
\text{PSNR}(x, \hat{x}) = 10 \cdot \log_{10} \left( \frac{\text{MAX}_I^2}{\text{MSE}(x, \hat{x})} \right)
\end{equation}
where $\text{MAX}_I = 1.0$ for floating-point normalized image representations.

The Structural Similarity Index (SSIM) evaluates perceptual fidelity across three orthogonal components---luminance $l(x, \hat{x})$, contrast $c(x, \hat{x})$, and structural correlation $s(x, \hat{x})$:
\begin{equation}
\text{SSIM}(x, \hat{x}) = \frac{(2\mu_x \mu_{\hat{x}} + C_1)(2\sigma_{x\hat{x}} + C_2)}{(\mu_x^2 + \mu_{\hat{x}}^2 + C_1)(\sigma_x^2 + \sigma_{\hat{x}}^2 + C_2)}
\end{equation}
where $\mu_x, \mu_{\hat{x}}$ denote local pixel mean intensities, $\sigma_x^2, \sigma_{\hat{x}}^2$ denote local variances, $\sigma_{x\hat{x}}$ denotes cross-covariance, and $C_1 = (0.01)^2, C_2 = (0.03)^2$.

When no pristine ground-truth reference image is available (as in real-world user uploads), noise standard deviation $\hat{\sigma}$ is estimated using the high-frequency Laplacian operator formulated by Immerkaer \cite{immerkaer1996fast}:
\begin{equation}
\hat{\sigma} = \sqrt{\frac{\pi}{2}} \cdot \frac{1}{6 (H - 2)(W - 2)} \sum_{i=2}^{H-1} \sum_{j=2}^{W-1} |(y * \mathcal{K}_{\text{lap}})(i, j)|
\end{equation}
where $\mathcal{K}_{\text{lap}}$ is the discrete $3 \times 3$ Laplacian kernel.

\section{Architecture and System Design}
\label{sec:architecture}

NeuroClean integrates three interconnected modules: (A) the Dual-Stream Hybrid U-Net--DnCNN neural network, (B) the 2D Hann-window overlapping cosine tiling engine, and (C) the full-stack asynchronous web application.

\subsection{Dual-Stream Hybrid Architecture}
The core neural engine combines two complementary feature extraction branches operating in parallel on the normalized input tensor $y \in \mathbb{R}^{B \times H_p \times W_p \times 3}$:

\subsubsection{Stream 1: Multi-Scale Contextual U-Net Branch}
The U-Net branch extracts hierarchical semantic and illumination context across contracting and expansive levels:
\begin{itemize}
    \item \textbf{Encoder Level 1:} Two $3 \times 3$ convolutions with 32 filters, Batch Normalization, and ReLU activations, followed by $2 \times 2$ max pooling ($H_p \to H_p/2$).
    \item \textbf{Encoder Level 2:} Two $3 \times 3$ convolutions with 64 filters, Batch Normalization, and ReLU activations, followed by $2 \times 2$ max pooling ($H_p/2 \to H_p/4$).
    \item \textbf{Bottleneck:} A $3 \times 3$ convolution with 128 filters, Batch Normalization, and ReLU activation.
    \item \textbf{Decoder Level 2:} A $2 \times 2$ nearest-neighbor upsampling layer ($H_p/4 \to H_p/2$), concatenated along the channel axis with Encoder Level 2 features via long skip connections, followed by two $3 \times 3$ convolutions with 64 filters, BN, and ReLU.
    \item \textbf{Decoder Level 1:} A $2 \times 2$ upsampling layer ($H_p/2 \to H_p$), concatenated with Encoder Level 1 features, followed by two $3 \times 3$ convolutions with 32 filters, producing contextual feature tensor $\Phi_{\text{unet}}(y) \in \mathbb{R}^{B \times H_p \times W_p \times 32}$.
\end{itemize}

\subsubsection{Stream 2: High-Frequency Residual DnCNN Branch}
The DnCNN branch preserves full spatial resolution to detect fine stochastic pixel grain:
\begin{itemize}
    \item \textbf{Layer 1:} A $3 \times 3$ convolution with 32 filters and ReLU activation (no Batch Normalization to preserve raw input distribution).
    \item \textbf{Layers 2--10:} Nine successive convolutional blocks, each comprising a $3 \times 3$ convolution with 32 filters, Batch Normalization, and ReLU activation.
    \item The output represents the fine-grained high-frequency feature tensor $\Phi_{\text{dncnn}}(y) \in \mathbb{R}^{B \times H_p \times W_p \times 32}$.
\end{itemize}

\subsubsection{Cross-Stream Feature Fusion and Residual Output}
The feature maps from both branches are concatenated at native spatial resolution:
\begin{equation}
\Phi_{\text{concat}}(y) = \left[ \Phi_{\text{unet}}(y) \, \parallel \, \Phi_{\text{dncnn}}(y) \right] \in \mathbb{R}^{B \times H_p \times W_p \times 64}
\end{equation}
A $1 \times 1$ cross-channel convolution reduces channel dimensionality while learning non-linear inter-branch dependencies:
\begin{equation}
\Phi_{\text{fused}}(y) = \text{ReLU}\left( \text{BN}\left( \text{Conv}_{1 \times 1}(\Phi_{\text{concat}}(y)) \right) \right) \in \mathbb{R}^{B \times H_p \times W_p \times 64}
\end{equation}
A final $3 \times 3$ refinement convolution projects the representation into the 3-channel estimated noise field:
\begin{equation}
\mathcal{R}_{\text{hybrid}}(y) = \text{Conv}_{3 \times 3}(\Phi_{\text{fused}}(y)) \in \mathbb{R}^{B \times H_p \times W_p \times 3}
\end{equation}
The restored patch estimate is obtained via residual subtraction:
\begin{equation}
\hat{x} = \text{clip}\left( y - \mathcal{R}_{\text{hybrid}}(y), \, 0.0, \, 1.0 \right)
\end{equation}

\subsection{2D Hann-Window Overlapping Cosine Tiling Engine}
To process high-resolution smartphone images of arbitrary dimensions $(H, W)$ without GPU memory overflow or tile boundary artifacts, we formulate Algorithm 1.

\begin{algorithm}[htbp]
\caption{2D Hann-Window Overlapping Cosine Tiling Inference}
\label{alg:tiling}
\begin{algorithmic}[1]
\REQUIRE Input image $y \in \mathbb{R}^{H \times W \times 3}$, tile size $T = 256$, step $S = 128$
\ENSURE Restored high-fidelity image $\hat{x} \in \mathbb{R}^{H \times W \times 3}$
\STATE Construct 1D Hann window: $w[i] = 0.5 \left(1 - \cos\left(\frac{2\pi (i + 0.5)}{T}\right)\right)$
\STATE $W_{\text{hann}} \leftarrow w \otimes w \in \mathbb{R}^{T \times T \times 1}$
\STATE Compute reflection padding $pad_h, pad_w$ for stride $S$
\STATE $y_{\text{pad}} \leftarrow \text{reflect\_pad}(y, pad_h, pad_w)$
\STATE Initialize accumulator buffers $\text{Output}_{\text{acc}} \leftarrow \mathbf{0}$, $\text{Weight}_{\text{acc}} \leftarrow \mathbf{0}$
\STATE Extract overlapping tiles $\{T_k\}$ at step $S$ with coordinates $\{(y_k, x_k)\}$
\STATE Run batched model prediction: $\{\hat{P}_k\} \leftarrow \text{Model.predict}(\{T_k\})$
\FOR{each tile prediction $\hat{P}_k$ at $(y_k, x_k)$}
    \STATE $\text{Output}_{\text{acc}}[y_k:y_k+T, x_k:x_k+T] += \hat{P}_k \odot W_{\text{hann}}$
    \STATE $\text{Weight}_{\text{acc}}[y_k:y_k+T, x_k:x_k+T] += W_{\text{hann}}$
\ENDFOR
\STATE Normalized image: $I_{\text{norm}} \leftarrow \text{Output}_{\text{acc}} / \max(\text{Weight}_{\text{acc}}, 10^{-6})$
\STATE $\hat{x} \leftarrow \text{clip}(I_{\text{norm}}[0:H, 0:W], 0.0, 1.0)$
\RETURN $\hat{x}$
\end{algorithmic}
\end{algorithm}

\section{Experimental Methodology}
\label{sec:methodology}

\subsection{Benchmark Dataset}
We evaluate our methodology on paired high-resolution smartphone photographs acquired under real-world low-light and high-ISO indoor/outdoor settings. The dataset comprises 100 paired clean and noisy smartphone image sets captured at native resolution ($1404 \times 2040$ pixels, 24-bit RGB):
\begin{itemize}
    \item \textbf{Clean References:} Captured at low ISO ($50 - 100$) mounted on a rigid tripod using long exposure times ($1.0 - 4.0\,\text{s}$).
    \item \textbf{Noisy Captures:} Captured under identical framing at high ISO ($1600 - 6400$) with fast shutter speeds ($1/50 - 1/200\,\text{s}$), inducing authentic photon shot noise, amplifier grain, and chroma mottling.
\end{itemize}

\subsection{Baseline Architectures for Comparison}
We benchmark five distinct configurations:
\begin{enumerate}
    \item \textbf{Noisy Input Baseline:} The uncleaned raw smartphone image.
    \item \textbf{Baseline 3-Layer CNN:} A sequential 3-layer convolutional network ($3 \times 3$ kernels, 64 filters, ReLUs) representing basic edge-filtering capacity (40,451 parameters, $7 \times 7$ receptive field).
    \item \textbf{Standard DnCNN (10 Layers):} A 10-layer feed-forward residual network with Batch Normalization (332,675 parameters, $21 \times 21$ receptive field).
    \item \textbf{Standard U-Net (3 Levels):} A 3-level contracting-expansive encoder-decoder with skip connections and spatial pooling (488,323 parameters, multi-scale receptive field).
    \item \textbf{Proposed Hybrid U-Net--DnCNN:} Our dual-stream architecture combining multi-scale context and high-frequency residual grain modeling (350,947 parameters, dual-scale global/local receptive field).
\end{enumerate}

\subsection{Training Protocol}
High-resolution training pairs are decomposed into $128 \times 128$ overlapping patches with a stride of 64 pixels. Data augmentation includes random horizontal flips, vertical flips, and $90^\circ$ orthogonal rotations. Models are trained using the Adam optimizer with an initial learning rate $\eta = 10^{-3}$, cosine decay schedule, batch size of 32, and Mean Squared Error (MSE) loss function for 50 epochs. Experiments were conducted on a Linux server equipped with an NVIDIA RTX 4090 GPU (24 GB VRAM) and Intel Core i9-13900K CPU @ 3.0 GHz with 64 GB DDR5 RAM.

\section{Experimental Results and Evaluation}
\label{sec:results}

\subsection{RQ1: Restoration Performance Comparison}
Table \ref{tab:benchmarks} summarizes the quantitative restoration metrics across the evaluated architectures.

\begin{table*}[htbp]
\centering
\caption{Quantitative restoration performance on smartphone benchmark dataset ($1404 \times 2040$ native resolution). Metrics: Mean PSNR (dB) and SSIM over 10 test evaluations.}
\label{tab:benchmarks}
\begin{tabular}{lcccccc}
\toprule
Architecture / Model & Parameters & Receptive Field & Mean PSNR (dB) & Mean SSIM & Latency (ms) & Status \\
\midrule
Noisy Baseline Input & --- & --- & $24.04 \pm 0.42$ & $0.7005$ & --- & Raw Capture \\
Baseline 3-Layer CNN & 40,451 & $7 \times 7$ & $26.87 \pm 0.31$ & $0.8441$ & $\sim 315$ ms & Comparative \\
Standard DnCNN (10 Layers) & 332,675 & $21 \times 21$ & $27.42 \pm 0.28$ & $0.8650$ & $\sim 480$ ms & Comparative \\
Standard U-Net (3 Levels) & 488,323 & Multi-Scale ($\sim 68$px) & $27.95 \pm 0.25$ & $0.8710$ & $\sim 520$ ms & Comparative \\
\textbf{Hybrid U-Net--DnCNN (Ours)} & \textbf{350,947} & \textbf{Dual-Scale (Global+Local)} & \textbf{28.60 $\pm$ 0.22} & \textbf{0.8920} & \textbf{$\sim 380$ ms} & \textbf{Proposed Model} \\
\bottomrule
\end{tabular}
\end{table*}

The proposed Hybrid U-Net--DnCNN achieves the highest restoration fidelity across all metrics:
\begin{itemize}
    \item \textbf{PSNR Gain:} Our hybrid architecture achieves \textbf{$28.60\,\text{dB}$}, outperforming the raw noisy baseline by $+4.56\,\text{dB}$, the 3-layer CNN by $+1.73\,\text{dB}$, standard DnCNN by $+1.18\,\text{dB}$, and standard U-Net by $+0.65\,\text{dB}$.
    \item \textbf{Structural Fidelity (SSIM):} Our model delivers \textbf{$0.8920$ SSIM}, representing a $+0.1915$ increase over the raw capture ($0.7005$) and surpassing standard U-Net ($0.8710$).
    \item \textbf{Statistical Significance:} Paired Wilcoxon signed-rank tests confirm that the PSNR and SSIM advantages of the hybrid model over all comparative baselines are statistically significant ($p_{\text{Holm}} < 0.001$, effect size $r > 0.85$).
\end{itemize}

\subsection{RQ2: Ablation Study on Architectural Components}
To isolate the contribution of each architectural innovation, we conduct an ablation study across four model variants (Table \ref{tab:ablation}).

\begin{table}[htbp]
\centering
\caption{Architectural ablation study on smartphone validation set.}
\label{tab:ablation}
\begin{tabular}{lcccc}
\toprule
Variant & Mean PSNR & Mean SSIM & Convergence & Params \\
\midrule
Variant A (U-Net Only) & $27.95$ dB & $0.8710$ & Epoch 34 & 245,603 \\
Variant B (DnCNN Only) & $27.42$ dB & $0.8650$ & Epoch 41 & 105,344 \\
Variant C (No 1x1 Fusion) & $28.18$ dB & $0.8785$ & Epoch 28 & 328,451 \\
Variant D (Direct Reconst.) & $27.65$ dB & $0.8682$ & Epoch 46 & 350,947 \\
\textbf{Full Hybrid Model} & \textbf{28.60 dB} & \textbf{0.8920} & \textbf{Epoch 22} & \textbf{350,947} \\
\bottomrule
\end{tabular}
\end{table}

\subsection{RQ3: Evaluation of Hann-Window Tiling vs. Naive Tiling}
We evaluate the performance of our 2D Hann-window overlapping cosine tiling engine against naive non-overlapping grid tiling across full-resolution $1404 \times 2040$ smartphone captures. We quantify boundary seam discontinuities using the Mean Boundary Gradient Discontinuity (MBGD) metric (Table \ref{tab:tiling}).

\begin{table}[htbp]
\centering
\caption{Comparison of tiling strategies on full-resolution smartphone captures ($1404 \times 2040$ pixels).}
\label{tab:tiling}
\begin{tabular}{lcccc}
\toprule
Strategy & Overlap & MBGD Seams & Mean PSNR & Peak VRAM \\
\midrule
No Tiling & --- & $0.0000$ & OOM Error & $>22.4$ GB \\
Naive Tiling & $0\%$ & $0.0842$ & $27.81$ dB & $410$ MB \\
Linear Blend & $50\%$ & $0.0185$ & $28.34$ dB & $580$ MB \\
\textbf{Hann Cosine} & \textbf{50\%} & \textbf{0.0012} & \textbf{28.60 dB} & \textbf{620 MB} \\
\bottomrule
\end{tabular}
\end{table}

\section{Discussion}
\label{sec:discussion}
\subsection{Implications for Practitioners}
\begin{enumerate}
    \item \textbf{Synergistic Feature Integration:} Coupling contextual contracting paths with full-resolution residual convolutional streams resolves macro-luminance splotches while safeguarding organic textures.
    \item \textbf{Decoupling Tile Inference from Border Artifacts:} In high-resolution mobile camera ISPs, Hann-window cosine weighting should be standard practice.
    \item \textbf{Production Web Microservices for ML Research:} Academic machine learning models should be accompanied by asynchronous web deployments to facilitate peer reproducibility.
\end{enumerate}

\section{Conclusion}
\label{sec:conclusion}
In this paper, we presented NeuroClean, a comprehensive deep learning framework and production web application for smartphone photograph denoising. On real-world paired smartphone photographs, NeuroClean achieves $28.60+\,\text{dB}$ PSNR and $0.8920+\,\text{SSIM}$, significantly outperforming classical and deep learning baselines ($p < 0.001$). Future work will focus on extending the hybrid architecture to RAW sensor Bayer-pattern inputs and quantizing the model to INT8 precision for on-device mobile ISP integration.

\section*{Data and Code Availability}
The software, trained checkpoints, and web deployment code are available on GitHub: \url{https://github.com/majjinaveen46/Image-Denoising-Smartphone}.

\begin{thebibliography}{00}
\bibitem{dudhane2022burst} A. Dudhane, S. W. Zamir, S. Khan, and F. S. Khan, ``Burst image restoration and enhancement,'' in \textit{Proc. IEEE/CVF Conf. Comput. Vis. Pattern Recognit. (CVPR)}, 2022. \url{https://arxiv.org/abs/2110.03680}
\bibitem{flepp2024real} R. Flepp, A. Ignatov, R. Timofte, and L. Van Gool, ``Real-world mobile image denoising dataset with efficient baselines,'' in \textit{Proc. IEEE/CVF Conf. Comput. Vis. Pattern Recognit. (CVPR)}, pp. 2596--2606, 2024. \url{https://paperswithcode.com/paper/real-world-mobile-image-denoising-dataset}
\bibitem{guo2024toward} S. Guo, ``Toward effective image restoration from raw domain,'' Ph.D. Dissertation, Dept. Comput., The Hong Kong Polytechnic Univ., 2024. \url{https://theses.lib.polyu.edu.hk/handle/200/12922}
\bibitem{li2023image} Z. Li, H. Liu, L. Cheng, and X. Jia, ``Image denoising algorithm based on gradient domain guided filtering and NSST,'' \textit{IEEE Access}, vol. 11, pp. 11923--11933, 2023. \url{https://ieeexplore.ieee.org/document/10035968}
\bibitem{kumari2024comparative} L. Kumari, N. Mittal, and M. Modi, ``Comparative analysis of noise reduction filters for medical images,'' in \textit{Proc. Signal Image Process. Appl.}, ResearchGate, 2024. \url{https://www.researchgate.net/publication/379162688_Comparative_Analysis_of_Noise_Reduction_Filters_for_Medical_Images}
\bibitem{lee2024development} J. Lee, ``Development of alternative filters to enhance the quality of digital images,'' \textit{ResearchGate / Digital Imaging Systems}, 2024. \url{https://www.researchgate.net/publication/379091587_Development_of_Alternative_Filters_to_Enhance_the_Quality_of_Digital_Images}
\bibitem{jebur2023review} R. S. Jebur, M. H. Bin Mohamed Zabil, D. A. Hammood, and L. K. Cheng, ``A comprehensive review of image denoising in deep learning,'' \textit{Multimedia Tools Appl.}, vol. 83, pp. 45671--45712, 2023. \url{https://link.springer.com/article/10.1007/s11042-023-17468-2}
\bibitem{mittal2024livenet} S. Mittal \textit{et al.}, ``LIVENet: A novel network for real-world low-light image denoising and enhancement,'' in \textit{Proc. IEEE/CVF Winter Conf. Appl. Comput. Vis. (WACV)}, pp. 2960--2970, 2024. \url{https://ieeexplore.ieee.org/document/10483686}
\bibitem{zhou2023deep} Z. Zhou, Y. Chen, and Y. Zhou, ``Deep dynamic memory augmented attentional dictionary learning for image denoising,'' \textit{IEEE Trans. Circuits Syst. Video Technol.}, vol. 33, no. 9, pp. 4668--4682, 2023. \url{https://ieeexplore.ieee.org/abstract/document/10054440}
\bibitem{gupta2023enhancing} O. Gupta, I. N. Kalapala, and K. M. Bhurchandi, ``Enhancing image denoising performance through a family of algorithms,'' in \textit{Proc. IEEE ISCON}, pp. 1--6, 2023. \url{https://ieeexplore.ieee.org/abstract/document/10136117}
\bibitem{zhou2023multiscale} L. Zhou, D. Zhou, H. Yang, S. Yang, and L. Nie, ``Multi-scale network toward real-world image denoising,'' \textit{Int. J. Mach. Learn. Cybern.}, vol. 14, pp. 1109--1123, 2023. \url{https://link.springer.com/article/10.1007/s13042-022-01694-5}
\bibitem{ulu2023mlfan} A. Ulu, G. Yildiz, and B. Dizdaroğlu, ``MLFAN: Multilevel feature attention network with texture prior for image denoising,'' \textit{IEEE Access}, vol. 11, pp. 35647--35661, 2023. \url{https://ieeexplore.ieee.org/abstract/document/10092851}
\bibitem{zhang2024application} G. Zhang, G. Zhang, Q. Lu, Y. Bo, and X. Zhou, ``Application of improved Canny algorithm in image edge detection of transmission lines,'' in \textit{Proc. IEEE ICICML}, pp. 312--316, 2024. \url{https://ieeexplore.ieee.org/abstract/document/10424911}
\bibitem{ranganath2023recurrent} A. Ranganath, F. Santiago, O. DeGuchy, M. Singhal, and R. F. Marcia, ``Recurrent neural imaging: An evolutionary approach for mixed Poisson-Gaussian image denoising,'' in \textit{Proc. IEEE DISCOVERY}, pp. 1--6, 2023. \url{https://ieeexplore.ieee.org/abstract/document/10069162/}
\bibitem{direk2024image} T. Direk, ``Image denoising by linear regression on non-local means algorithm,'' \textit{Signal, Image Video Process.}, vol. 18, pp. 3185--3194, 2024. \url{https://link.springer.com/article/10.1007/s11760-024-03086-4}
\bibitem{delbracio2021mobile} M. Delbracio, D. Kelly, M. S. Brown, and P. Milanfar, ``Mobile computational photography: A tour,'' \textit{Annu. Rev. Vis. Sci.}, vol. 7, pp. 571--604, 2021.
\bibitem{zhang2017beyond} K. Zhang, W. Zuo, Y. Chen, D. Meng, and L. Zhang, ``Beyond a Gaussian denoiser: Residual learning of deep CNN for image denoising,'' \textit{IEEE Trans. Image Process.}, vol. 26, no. 7, pp. 3142--3155, 2017.
\bibitem{ronneberger2015u} O. Ronneberger, P. Fischer, and T. Brox, ``U-Net: Convolutional networks for biomedical image segmentation,'' in \textit{Proc. MICCAI}, Springer, 2015, pp. 234--241.
\bibitem{foi2008practical} A. Foi, M. Trimeche, V. Katkovnik, and K. Egiazarian, ``Practical Poissonian-Gaussian noise modeling and fitting for single-image raw-data,'' \textit{IEEE Trans. Image Process.}, vol. 17, no. 10, pp. 1737--1754, 2008.
\bibitem{abdelhamed2018high} A. Abdelhamed, S. Lin, and M. S. Brown, ``A high-quality smartphone dataset for image denoising,'' in \textit{Proc. CVPR}, 2018, pp. 1692--1700.
\bibitem{huang2021evaluation} Y. Huang \textit{et al.}, ``Evaluation of tiling strategies for deep learning image restoration on ultra-high-resolution images,'' \textit{IEEE Access}, vol. 9, pp. 84120--84131, 2021.
\bibitem{immerkaer1996fast} J. Immerkaer, ``Fast noise variance estimation,'' \textit{Comput. Vis. Image Underst.}, vol. 64, no. 2, pp. 300--302, 1996.
\end{thebibliography}

\begin{IEEEbiography}[{\includegraphics[width=1in,height=1.25in,clip,keepaspectratio]{author1.png}}]{Majji Naveen} is currently pursuing a Bachelor of Technology degree in Computer Science and Engineering at Vellore Institute of Technology (VIT), Vellore, India (Email: majji.naveen2023@gmail.com). His research interests encompass deep learning, computational photography, computer vision, and high-performance machine learning system deployment. Through his work on the NeuroClean project, he has pioneered hybrid neural architectures combining multi-scale contextual representations with deep residual noise estimation, alongside engineered windowed inference algorithms for high-resolution computational imaging.
\end{IEEEbiography}

\begin{IEEEbiography}[{\includegraphics[width=1in,height=1.25in,clip,keepaspectratio]{author2.png}}]{Parri Hari Chandan} is currently pursuing a B.Tech. degree in Computer Science and Engineering at Vellore Institute of Technology (VIT), Vellore, India (Email: harichandan.parri2023@vitstudent.ac.in). Research interests include digital signal processing, convolutional neural networks, and web-based artificial intelligence deployment frameworks.
\end{IEEEbiography}

\begin{IEEEbiography}[{\includegraphics[width=1in,height=1.25in,clip,keepaspectratio]{author3.png}}]{M. Siva Charan Reddy} is currently pursuing a B.Tech. degree in Computer Science and Engineering at Vellore Institute of Technology (VIT), Vellore, India (Email: mareddy.siva2023@vitstudent.ac.in). Research interests include computer vision, neural image enhancement, and cloud-based machine learning systems.
\end{IEEEbiography}

\begin{IEEEbiography}[{\includegraphics[width=1in,height=1.25in,clip,keepaspectratio]{author4.png}}]{Jayalakshmi M} is an Associate Professor at the School of Computer Science and Engineering (SCOPE), Vellore Institute of Technology (VIT), Vellore, India (Email: m.jayalakshmi@vit.ac.in). With extensive academic and research experience in artificial intelligence, computer vision, digital image processing, and distributed machine learning systems, she has supervised numerous research initiatives and published extensively in indexed international journals and conferences.
\end{IEEEbiography}

\EOD

\end{document}
